% ======================================================================
%  UPPGIFT 3: EULERCYKEL
% ======================================================================
\section{3 Eulercykel}

\uppgiftsbild{3}{Eulercykel}
{hitta en Eulercykel i linjär tid — och kunna säga vad ”linjär” betyder här.}
{\begin{itemize}\setlength{\itemsep}{1pt}
  \item Varje \emph{kant} exakt en gång. (Varje \emph{hörn} en gång är ett helt annat problem.)
  \item Varför jämnt gradtal räcker: man kan bara fastna där man började.
  \item Linjärt i $|V|+|E|$ kräver rätt datastruktur: delade kanter och en pekare per hörn.
\end{itemize}}

\begin{frame}{Uppgift 3: Eulercykel}
\begin{infobox}
Givet en oriktad sammanhängande graf $G=\langle V,E\rangle$ där alla hörn har jämnt gradtal.
Hitta en Eulercykel, dvs.\ en sluten stig som passerar varje kant i $E$ exakt en gång.
Algoritmen ska gå i linjär tid.
\end{infobox}
\begin{columns}[T]
\begin{column}{0.5\textwidth}
\centering
\begin{tikzpicture}[scale=0.85]
  \node[gv] (1) at (0,0) {1};
  \node[gv] (2) at (1.5,1.2) {2};
  \node[gv] (3) at (3,0) {3};
  \node[gv] (4) at (-1.5,1.2) {4};
  \node[gv] (5) at (-1.5,-1.2) {5};
  \node[gv] (6) at (4.5,1.2) {6};
  \node[gv] (7) at (4.5,-1.2) {7};
  \draw[ge] (1) -- (2);
  \draw[ge] (2) -- (3);
  \draw[ge] (3) -- (1);
  \draw[ge] (1) -- (4);
  \draw[ge] (4) -- (5);
  \draw[ge] (5) -- (1);
  \draw[ge] (3) -- (6);
  \draw[ge] (6) -- (7);
  \draw[ge] (7) -- (3);
  \node[font=\scriptsize,text=kthorange] at (0,-0.5) {4};
  \node[font=\scriptsize,text=kthorange] at (1.5,1.7) {2};
  \node[font=\scriptsize,text=kthorange] at (3,-0.5) {4};
  \node[font=\scriptsize,text=kthorange] at (-2.0,1.2) {2};
  \node[font=\scriptsize,text=kthorange] at (-2.0,-1.2) {2};
  \node[font=\scriptsize,text=kthorange] at (5.0,1.2) {2};
  \node[font=\scriptsize,text=kthorange] at (5.0,-1.2) {2};
\end{tikzpicture}
\par
{\scriptsize\color{kthorange}gradtal i orange}
\end{column}
\begin{column}{0.47\textwidth}
\small
\begin{itemize}\setlength{\itemsep}{3pt}
  \item Hörn får upprepas: här måste 1 och 3 passeras två gånger var.
  \item Villkoren är nödvändiga: varje gång cykeln går \emph{in} i ett hörn går den också
        \emph{ut}, så gradtalen blir jämna.
  \item Att de också är tillräckliga visar algoritmen.
\end{itemize}
\end{column}
\end{columns}
\end{frame}

\begin{frame}{Euler eller Hamilton? Ett ord skiljer}
\begin{columns}[T]
\begin{column}{0.3\textwidth}
\centering
{\small\bfseries\color{kthnavy}Fluga}\par\smallskip
\begin{tikzpicture}[scale=0.85]
  \node[gv] (m) at (0,0) {}; \node[gv] (a) at (-1.3,0.85) {}; \node[gv] (b) at (-1.3,-0.85) {};
  \node[gv] (c) at (1.3,0.85) {}; \node[gv] (d) at (1.3,-0.85) {};
  \draw[ge] (m)--(a)--(b)--(m)--(c)--(d)--(m);
\end{tikzpicture}\par\smallskip
\footnotesize Euler: \ja\quad Hamilton: \nej\\ Mitten måste passeras två gånger.
\end{column}
\begin{column}{0.3\textwidth}
\centering
{\small\bfseries\color{kthnavy}$K_4$}\par\smallskip
\begin{tikzpicture}[scale=0.85]
  \node[gv] (a) at (0,0) {}; \node[gv] (b) at (1.7,0) {};
  \node[gv] (c) at (1.7,1.7) {}; \node[gv] (d) at (0,1.7) {};
  \draw[ge] (a)--(b)--(c)--(d)--(a)--(c) (b)--(d);
\end{tikzpicture}\par\smallskip
\footnotesize Euler: \nej\quad Hamilton: \ja\\ Alla gradtal är 3.
\end{column}
\begin{column}{0.36\textwidth}
\small
\textbf{Euler:} varje \emph{kant} exakt en gång. Linjär tid.\par\smallskip
\textbf{Hamilton:} varje \emph{hörn} exakt en gång. NP-fullständigt.\par\medskip
Nästan samma mening. Helt olika problem.
\end{column}
\end{columns}
\vfill
\begin{missbox}
”En Eulercykel besöker alla hörn en gång.” Det är en Hamiltoncykel — och för den finns
ingen känd polynomisk algoritm.
\end{missbox}
\varstrip{Kontrast}{”besök allt exakt en gång”}{kanter eller hörn}{vad ordet ”varje” syftar på}
\end{frame}

\begin{frame}{Idén: en hare och en sköldpadda}
\begin{columns}[T]
\begin{column}{0.49\textwidth}
\begin{infobox}[title={\textbf{Haren} (Viggos \textsc{PathFinder})},coltitle=white,colbacktitle=kthorange,
  colframe=kthorange,fonttitle=\small,toptitle=1pt,bottomtitle=1pt]
Springer från ett hörn $u$ längs omärkta kanter, märker dem och stannar först när den är
tillbaka i $u$. Jämna gradtal gör att den bara kan fastna där.
\end{infobox}
\end{column}
\begin{column}{0.49\textwidth}
\begin{infobox}[title={\textbf{Sköldpaddan} (Viggos \textsc{Straggler})},coltitle=white,colbacktitle=kthgreen,
  colframe=kthgreen,fonttitle=\small,toptitle=1pt,bottomtitle=1pt]
Går efter längs harens stig och skriver ut hörnen. Har ett hörn omärkta kanter kvar skickar
den ut haren därifrån — och går den nya turen \emph{innan} den fortsätter.
\end{infobox}
\end{column}
\end{columns}
\medskip
\small
Varje tur är en sluten stig. Sköldpaddan skarvar in turerna i varandra, precis där de börjar.
Ingen kant springs två gånger av haren, och ingen kant gås två gånger av sköldpaddan.
\medskip
\begin{ahabox}
Därför linjärt: varje kant märks en gång av haren och passeras en gång av sköldpaddan.
(Förutsatt att båda hittar nästa omärkta kant i konstant tid — mer om det strax.)
\end{ahabox}
\end{frame}

\begin{frame}{Pseudokod}
\begin{columns}[T]
\begin{column}{0.5\textwidth}
{\small\bfseries\color{kthnavy}EulerCykel$(G)$}
\footnotesize
\begin{algorithmic}[1]
\State $\mathit{cykel}\gets[v_0]$
\State välj en kant $(v_0,t)$ och märk den
\State \textsc{Sköldpadda}(\textsc{Hare}$(v_0,t)$)
\State \Return $\mathit{cykel}$
\end{algorithmic}
\medskip
{\small\bfseries\color{kthorange}Hare$(\mathit{start},\mathit{cur})$}
\footnotesize
\begin{algorithmic}[1]
\State $\mathit{stig}\gets[\mathit{cur}]$
\While{$\mathit{cur}\neq\mathit{start}$}
  \State välj en omärkt kant $(\mathit{cur},v)$; märk den
  \State lägg $v$ sist i $\mathit{stig}$;\ $\mathit{cur}\gets v$
\EndWhile
\State \Return $\mathit{stig}$
\end{algorithmic}
\end{column}
\begin{column}{0.47\textwidth}
{\small\bfseries\color{kthgreen}Sköldpadda$(\mathit{stig})$}
\footnotesize
\begin{algorithmic}[1]
\ForAll{$u$ i $\mathit{stig}$, i ordning}
  \State lägg $u$ sist i $\mathit{cykel}$
  \ForAll{kanter $(u,v)$}
    \If{$(u,v)$ är omärkt}
      \State märk $(u,v)$
      \State \textsc{Sköldpadda}(\textsc{Hare}$(u,v)$)
    \EndIf
  \EndFor
\EndFor
\end{algorithmic}
\medskip
\begin{infobox}[fontupper=\footnotesize]
Omärkt-testet görs om för varje kant: det rekursiva anropet kan ha märkt kanter som slingan
ännu inte kommit till.
\end{infobox}
\end{column}
\end{columns}
\end{frame}

\begin{frame}{Eulercykeln steg för steg}
\begin{columns}[T]
\begin{column}{0.54\textwidth}
\centering
\begin{tikzpicture}[scale=0.92]
  \useasboundingbox (-2.1,-1.75) rectangle (5.1,1.75);
  \node[gv] (1) at (0,0) {1};
  \node[gv] (2) at (1.5,1.2) {2};
  \node[gv] (3) at (3,0) {3};
  \node[gv] (4) at (-1.5,1.2) {4};
  \node[gv] (5) at (-1.5,-1.2) {5};
  \node[gv] (6) at (4.5,1.2) {6};
  \node[gv] (7) at (4.5,-1.2) {7};
  \draw[faint] (1) -- (2);
  \draw[faint] (2) -- (3);
  \draw[faint] (3) -- (1);
  \draw[faint] (1) -- (4);
  \draw[faint] (4) -- (5);
  \draw[faint] (5) -- (1);
  \draw[faint] (3) -- (6);
  \draw[faint] (6) -- (7);
  \draw[faint] (7) -- (3);
  \draw[->,draw=kthblue,line width=2.2pt,shorten >=1pt,visible on=<1->] (1) -- (2);
  \draw[->,draw=kthblue,line width=3.6pt,opacity=0.35,visible on=<1>] (1) -- (2);
  \draw[->,draw=kthblue,line width=2.2pt,shorten >=1pt,visible on=<1->] (2) -- (3);
  \draw[->,draw=kthblue,line width=3.6pt,opacity=0.35,visible on=<1>] (2) -- (3);
  \draw[->,draw=kthblue,line width=2.2pt,shorten >=1pt,visible on=<1->] (3) -- (1);
  \draw[->,draw=kthblue,line width=3.6pt,opacity=0.35,visible on=<1>] (3) -- (1);
  \draw[->,draw=kthgrass,line width=2.2pt,shorten >=1pt,visible on=<3->] (3) -- (6);
  \draw[->,draw=kthgrass,line width=3.6pt,opacity=0.35,visible on=<3>] (3) -- (6);
  \draw[->,draw=kthgrass,line width=2.2pt,shorten >=1pt,visible on=<3->] (6) -- (7);
  \draw[->,draw=kthgrass,line width=3.6pt,opacity=0.35,visible on=<3>] (6) -- (7);
  \draw[->,draw=kthgrass,line width=2.2pt,shorten >=1pt,visible on=<3->] (7) -- (3);
  \draw[->,draw=kthgrass,line width=3.6pt,opacity=0.35,visible on=<3>] (7) -- (3);
  \draw[->,draw=kthorange,line width=2.2pt,shorten >=1pt,visible on=<5->] (1) -- (4);
  \draw[->,draw=kthorange,line width=3.6pt,opacity=0.35,visible on=<5>] (1) -- (4);
  \draw[->,draw=kthorange,line width=2.2pt,shorten >=1pt,visible on=<5->] (4) -- (5);
  \draw[->,draw=kthorange,line width=3.6pt,opacity=0.35,visible on=<5>] (4) -- (5);
  \draw[->,draw=kthorange,line width=2.2pt,shorten >=1pt,visible on=<5->] (5) -- (1);
  \draw[->,draw=kthorange,line width=3.6pt,opacity=0.35,visible on=<5>] (5) -- (1);
  \node[circle,draw=kthgreen,line width=1.6pt,minimum size=8.4mm,visible on=<2>] at (2) {};
  \node[circle,draw=kthgreen,line width=1.6pt,minimum size=8.4mm,visible on=<3>] at (3) {};
  \node[circle,draw=kthgreen,line width=1.6pt,minimum size=8.4mm,visible on=<4>] at (3) {};
  \node[circle,draw=kthgreen,line width=1.6pt,minimum size=8.4mm,visible on=<5>] at (1) {};
  \node[circle,draw=kthgreen,line width=1.6pt,minimum size=8.4mm,visible on=<6>] at (1) {};
  \node[circle,fill=white,draw=black!50,inner sep=0.8pt,font=\scriptsize\bfseries,visible on=<6>] at ($(1)!0.5!(2)$) {1};
  \node[circle,fill=white,draw=black!50,inner sep=0.8pt,font=\scriptsize\bfseries,visible on=<6>] at ($(2)!0.5!(3)$) {2};
  \node[circle,fill=white,draw=black!50,inner sep=0.8pt,font=\scriptsize\bfseries,visible on=<6>] at ($(3)!0.5!(1)$) {6};
  \node[circle,fill=white,draw=black!50,inner sep=0.8pt,font=\scriptsize\bfseries,visible on=<6>] at ($(1)!0.5!(4)$) {7};
  \node[circle,fill=white,draw=black!50,inner sep=0.8pt,font=\scriptsize\bfseries,visible on=<6>] at ($(4)!0.5!(5)$) {8};
  \node[circle,fill=white,draw=black!50,inner sep=0.8pt,font=\scriptsize\bfseries,visible on=<6>] at ($(5)!0.5!(1)$) {9};
  \node[circle,fill=white,draw=black!50,inner sep=0.8pt,font=\scriptsize\bfseries,visible on=<6>] at ($(3)!0.5!(6)$) {3};
  \node[circle,fill=white,draw=black!50,inner sep=0.8pt,font=\scriptsize\bfseries,visible on=<6>] at ($(6)!0.5!(7)$) {4};
  \node[circle,fill=white,draw=black!50,inner sep=0.8pt,font=\scriptsize\bfseries,visible on=<6>] at ($(7)!0.5!(3)$) {5};
\end{tikzpicture}
\par\smallskip{\scriptsize\textcolor{kthblue}{\textbf{tur 1}}\quad\textcolor{kthgrass}{\textbf{tur 2}}\quad\textcolor{kthorange}{\textbf{tur 3}}\quad\textcolor{kthgreen}{$\bigcirc$ sköldpaddan}}
\end{column}
\begin{column}{0.43\textwidth}
\only<1>{%
{\small\textbf{Senaste tur (haren):}} $[2,3,1]$\par\medskip
{\small\textbf{Cykel (sköldpaddan):}}\par
$\mathit{cykel}=[1]$\par\medskip
\begin{infobox}[fontupper=\footnotesize]Märk $(1,2)$. \textbf{Haren} springer längs omärkta kanter: $1\to2\to3\to1$. Den fastnar i 1 — starthörnet.\end{infobox}}
\only<2>{%
{\small\textbf{Senaste tur (haren):}} $[2,3,1]$\par\medskip
{\small\textbf{Cykel (sköldpaddan):}}\par
$\mathit{cykel}=[1,\textcolor{kthorange}{\mathbf{2}}]$\par\medskip
\begin{infobox}[fontupper=\footnotesize]\textbf{Sköldpaddan} går efter längs stigen. I 2: alla kanter märkta. Gå vidare.\end{infobox}}
\only<3>{%
{\small\textbf{Senaste tur (haren):}} $[6,7,3]$\par\medskip
{\small\textbf{Cykel (sköldpaddan):}}\par
$\mathit{cykel}=[1,2,\textcolor{kthorange}{\mathbf{3}}]$\par\medskip
\begin{infobox}[fontupper=\footnotesize]I 3: omärkta kanter kvar! Skicka ut haren från 3: $3\to6\to7\to3$.\end{infobox}}
\only<4>{%
{\small\textbf{Senaste tur (haren):}} $[6,7,3]$\par\medskip
{\small\textbf{Cykel (sköldpaddan):}}\par
$\mathit{cykel}=[1,2,3,\textcolor{kthorange}{\mathbf{6}},\textcolor{kthorange}{\mathbf{7}},\textcolor{kthorange}{\mathbf{3}}]$\par\medskip
\begin{infobox}[fontupper=\footnotesize]Sköldpaddan går den nya turen \emph{innan} den fortsätter: 6, 7, 3.\end{infobox}}
\only<5>{%
{\small\textbf{Senaste tur (haren):}} $[4,5,1]$\par\medskip
{\small\textbf{Cykel (sköldpaddan):}}\par
$\mathit{cykel}=[1,2,3,6,7,3,\textcolor{kthorange}{\mathbf{1}}]$\par\medskip
\begin{infobox}[fontupper=\footnotesize]Tillbaka på första turen. I 1: omärkta kanter. Haren: $1\to4\to5\to1$.\end{infobox}}
\only<6>{%
{\small\textbf{Senaste tur (haren):}} $[4,5,1]$\par\medskip
{\small\textbf{Cykel (sköldpaddan):}}\par
$\mathit{cykel}=[1,2,3,6,7,3,1,\textcolor{kthorange}{\mathbf{4}},\textcolor{kthorange}{\mathbf{5}},\textcolor{kthorange}{\mathbf{1}}]$\par\medskip
\begin{infobox}[fontupper=\footnotesize]Sköldpaddan går 4, 5, 1. Klart: $10=|E|+1$ hörn. Siffrorna visar cykelns ordning.\end{infobox}}
\end{column}
\end{columns}
\varstrip{Generalisering}{algoritmen}{var turerna börjar}{att turer skarvas in där sköldpaddan hittar omärkta kanter}
\end{frame}


{\kaninbg
\begin{frame}{\kh Varför fastnar haren bara i starthörnet?}
\begin{columns}[T]
\begin{column}{0.55\textwidth}
\small
\begin{itemize}\setlength{\itemsep}{3pt}
  \item Varje gång haren kommer \emph{in} i ett hörn $w\neq u$ har den använt ett udda antal
        av $w$:s kanter (in, ut, in, ut, \dots, in).
  \item $\deg(w)$ är jämnt, så minst en kant är omärkt. Haren kommer alltid vidare.
  \item I $u$ är det tvärtom: den första kanten gick \emph{ut}, så efter varje ankomst är ett
        jämnt antal använt. Där kan den fastna.
  \item En sluten tur tar bort ett jämnt antal kanter från varje hörn. Alla gradtal är
        fortfarande jämna, och argumentet gäller nästa tur också.
\end{itemize}
\end{column}
\begin{column}{0.42\textwidth}
\begin{kaninbox}[fontupper=\footnotesize]
\textbf{Och varför kommer alla kanter med?} Anta att en kant blev omärkt. Grafen är
sammanhängande, så någon omärkt kant sitter i ett hörn som finns i cykeln. Men sköldpaddan
kontrollerade varje hörn i cykeln. Motsägelse.
\end{kaninbox}
\begin{kaninbox}[fontupper=\footnotesize]
\textbf{Udda gradtal:} har exakt två hörn udda grad finns en Eulerväg mellan dem. Lägg till
en extra kant mellan dem, hitta en cykel, ta bort kanten.
\end{kaninbox}
\end{column}
\end{columns}
\end{frame}}

\begin{frame}{Linjär i vad? Tre detaljer som avgör}
\begin{columns}[T]
\begin{column}{0.6\textwidth}
\footnotesize
\textbf{Indata:} grannlistor, $|V|+|E|$. \textbf{Utdata:} $|E|+1$ hörn, så $\Omega(|E|)$ krävs ändå.
\begin{enumerate}\setlength{\itemsep}{2pt}
  \item \textbf{En kant, två listor.} $\{u,v\}$ står i både $u$:s och $v$:s lista. Lagra kanten
        en gång, med markeringen; båda listorna pekar på den.
  \item \textbf{En pekare per hörn} till nästa ej undersökta plats i listan. Utan den börjar
        haren och sköldpaddan om från början varje gång.
  \item Med båda: varje listplats passeras ett konstant antal gånger. $O(|V|+|E|)$.
\end{enumerate}
\end{column}
\begin{column}{0.37\textwidth}
\centering
\begin{tikzpicture}[scale=0.62]
  \node[gv,fill=lampyellow!50] (c) at (0,0) {$c$};
  \foreach \a [count=\i] in {45,135,225,315} {
    \node[gv,minimum size=3.5mm] (x\i) at (\a-22:1.7) {};
    \node[gv,minimum size=3.5mm] (y\i) at (\a+22:1.7) {};
    \draw[ge] (c)--(x\i)--(y\i)--(c);
  }
\end{tikzpicture}\par
\footnotesize Blomma med $k$ blad: $\deg(c)=2k$, och cykeln passerar $c$ $k$ gånger.
Utan pekare: $k\cdot2k=\Theta(|E|^2)$.
\end{column}
\end{columns}
\vfill
\begin{missbox}[fontupper=\footnotesize]
”Eulercykeln hittas i $O(n)$.” Med $n=|V|$ är det omöjligt: $K_{101}$ har 101 hörn och
5\,050 kanter, och alla ska med i cykeln.
\end{missbox}
\varstrip{Separation}{algoritmen}{datastrukturen: delade kanter, pekare}{att ”linjär” kräver rätt representation}
\end{frame}

{\kaninbg
\begin{frame}{\kh Utan rekursion: Hierholzer med stack}
\begin{columns}[T]
\begin{column}{0.52\textwidth}
{\small\bfseries\color{kthnavy}Hierholzer$(G,v_0)$}
\footnotesize
\begin{algorithmic}[1]
\State $S\gets[v_0]$;\ $\mathit{cykel}\gets[\,]$
\While{$S\neq\emptyset$}
  \State $u\gets$ översta elementet i $S$
  \If{$u$ har en omärkt kant $(u,v)$} \Comment{pekaren}
    \State märk $(u,v)$;\ lägg $v$ överst i $S$
  \Else
    \State ta bort $u$ ur $S$;\ lägg $u$ sist i $\mathit{cykel}$
  \EndIf
\EndWhile
\State \Return $\mathit{cykel}$
\end{algorithmic}
\end{column}
\begin{column}{0.45\textwidth}
\small
Stacken är harens stig. När haren fastnar backar vi och skriver ut hörnen.\par\medskip
På exemplet: $[1,5,4,1,3,7,6,3,2,1]$ — samma cykel som förut, baklänges. En Eulercykel
baklänges är också en Eulercykel.
\end{column}
\end{columns}
\vfill
\begin{kaninbox}[fontupper=\footnotesize]
\textbf{Varför bry sig?} Rekursionen i hare/sköldpadda kan bli $\Theta(|E|)$ djup. Med en
miljon kanter tar anropsstacken slut i Java och Python långt innan tiden gör det.
\end{kaninbox}
\begin{kaninbox}[fontupper=\footnotesize]
\textbf{Kontrast: Fleury (1883).} Gå aldrig över en bro (en kant vars borttagning delar grafen)
om det finns ett annat val. Korrekt, men varje steg kräver en brokontroll: $O(|E|^2)$.
\end{kaninbox}
\end{frame}}

{\kaninbg
\begin{frame}{\kh Königsberg 1736: där allt började}
\begin{columns}[T]
\begin{column}{0.42\textwidth}
\centering
\begin{tikzpicture}[scale=0.9]
  \node[gv] (A) at (0,0) {$A$};   \node[gv] (B) at (0,1.9) {$B$};
  \node[gv] (C) at (0,-1.9) {$C$}; \node[gv] (D) at (2.7,0) {$D$};
  \draw[ge] (A) to[bend left=22] (B);  \draw[ge] (A) to[bend right=22] (B);
  \draw[ge] (A) to[bend left=22] (C);  \draw[ge] (A) to[bend right=22] (C);
  \draw[ge] (A) -- (D); \draw[ge] (B) -- (D); \draw[ge] (C) -- (D);
  \foreach \v/\d/\p in {A/5/left,B/3/left,C/3/left,D/3/right}
    \node[font=\scriptsize,text=kthorange,\p=4pt of \v] {\d};
\end{tikzpicture}
\end{column}
\begin{column}{0.55\textwidth}
\small
\begin{itemize}\setlength{\itemsep}{3pt}
  \item Fyra landområden, sju broar. Går det att promenera över varje bro exakt en gång?
  \item Euler: nej. Alla fyra hörnen har udda grad (5, 3, 3, 3). För en cykel får inget ha
        udda grad, för en väg högst två.
  \item Eulers observation räknas som grafteorins första sats. Han bevisade bara att villkoret
        är nödvändigt. Att det räcker visades senare (Hierholzer, 1873).
\end{itemize}
\end{column}
\end{columns}
\vfill
\begin{kaninbox}[fontupper=\footnotesize]
$A$ är ön Kneiphof, $B$ och $C$ flodstränderna, $D$ området öster om ön. Euler behövde aldrig
promenera: det räckte att räkna gradtal.
\end{kaninbox}
\end{frame}}

{\kaninbg
\begin{frame}{\kh Kodlåset: Eulercykler i riktade grafer}
\begin{columns}[T]
\begin{column}{0.45\textwidth}
\centering
\begin{tikzpicture}[scale=0.85,every node/.style={font=\small},
  lab/.style={font=\scriptsize,fill=white,inner sep=0.8pt}]
  \node[gv,minimum size=7mm] (00) at (0,0) {00};
  \node[gv,minimum size=7mm] (01) at (1.8,1.1) {01};
  \node[gv,minimum size=7mm] (10) at (1.8,-1.1) {10};
  \node[gv,minimum size=7mm] (11) at (3.6,0) {11};
  \draw[->,ge] (00) to[out=150,in=210,looseness=5] node[lab,left] {000} (00);
  \draw[->,ge] (00) -- node[lab] {001} (01);
  \draw[->,ge] (01) to[bend left=18] node[lab,right] {010} (10);
  \draw[->,ge] (10) to[bend left=18] node[lab,left] {101} (01);
  \draw[->,ge] (01) -- node[lab] {011} (11);
  \draw[->,ge] (11) to[out=30,in=-30,looseness=5] node[lab,right] {111} (11);
  \draw[->,ge] (11) -- node[lab] {110} (10);
  \draw[->,ge] (10) -- node[lab] {100} (00);
\end{tikzpicture}
\end{column}
\begin{column}{0.52\textwidth}
\small
\begin{itemize}\setlength{\itemsep}{2pt}
  \item Ett kodlås öppnas så fort de fyra senaste tryckningarna är rätt. Hur många tryck
        krävs för att pröva alla $10^4$ koder?
  \item Hörn: tre siffror. Kant: en fjärde. Ingrad = utgrad = 10 överallt, så en riktad
        Eulercykel finns.
  \item Svar: $10^4+3=10\,003$ tryck i stället för $40\,000$.
  \item Binärt, tre siffror (bilden): $0001011100$ innehåller alla åtta.
\end{itemize}
\end{column}
\end{columns}
\vfill
\begin{kaninbox}[fontupper=\footnotesize]
Samma algoritm, riktade grafer: de Bruijn-följder. (Kodlås borde alltså kräva Enter.)
\end{kaninbox}
\varstrip{Generalisering}{algoritmen}{grafen: oriktad eller riktad}{att villkoret är ”lika mycket in som ut”}
\end{frame}}

\begin{frame}{Båda villkoren behövs}
\begin{columns}[T]
\begin{column}{0.31\textwidth}
\centering
{\footnotesize\bfseries\color{kthnavy}Jämna grader, två delar}\par\smallskip
\begin{tikzpicture}[scale=0.9]
  \node[gv] (a) at (0,0) {}; \node[gv] (b) at (1.1,0) {}; \node[gv] (c) at (0.55,0.95) {};
  \node[gv] (d) at (1.9,0) {}; \node[gv] (e) at (3.0,0) {}; \node[gv] (f) at (2.45,0.95) {};
  \draw[ge] (a)--(b)--(c)--(a) (d)--(e)--(f)--(d);
\end{tikzpicture}\par\medskip
\footnotesize Ingen Eulercykel.
\end{column}
\begin{column}{0.31\textwidth}
\centering
{\footnotesize\bfseries\color{kthnavy}Sammanhängande, två udda}\par\smallskip
\begin{tikzpicture}[scale=0.8]
  \node[gv,fill=coral!40] (A) at (0,0) {}; \node[gv,fill=coral!40] (B) at (1.6,0) {};
  \node[gv] (C) at (0,1.5) {}; \node[gv] (D) at (1.6,1.5) {}; \node[gv] (E) at (0.8,2.4) {};
  \draw[ge] (A)--(B)--(D)--(C)--(A)--(D) (B)--(C) (C)--(E)--(D);
\end{tikzpicture}\par
\footnotesize Nikolaus hus: Eulerväg mellan de röda, men ingen cykel.
\end{column}
\begin{column}{0.31\textwidth}
\centering
{\footnotesize\bfseries\color{kthnavy}Sammanhängande, jämna}\par\smallskip
\begin{tikzpicture}[scale=0.8]
  \node[gv] (m) at (0,0) {}; \node[gv] (a) at (-1.2,0.8) {}; \node[gv] (b) at (-1.2,-0.8) {};
  \node[gv] (c) at (1.2,0.8) {}; \node[gv] (d) at (1.2,-0.8) {};
  \draw[ge] (m)--(a)--(b)--(m)--(c)--(d)--(m);
\end{tikzpicture}\par\medskip
\footnotesize Eulercykel finns.
\end{column}
\end{columns}
\vfill
\begin{missbox}
”Alla grader jämna räcker.” Två trianglar utan förbindelse har jämna grader och ingen
Eulercykel. (Isolerade hörn gör däremot inget: de har inga kanter att passera.)
\end{missbox}
\varstrip{Separation}{frågan}{ett villkor i taget}{att både sammanhang och jämna grader behövs}
\end{frame}

\begin{frame}{Vanliga fel i uppgift 3}
\begin{columns}[T]
\begin{column}{0.49\textwidth}
\begin{missbox}[fontupper=\footnotesize]
”Gå tills du fastnar, klart.” Första turen $1\to2\to3\to1$ tar 3 av 9 kanter. Man fastnar i
starthörnet, inte när grafen är slut.
\end{missbox}
\begin{missbox}[fontupper=\footnotesize]
”Märk kanten i $u$:s lista.” Den står i $v$:s lista också. Märks bara ena riktningen går
kanten två gånger.
\end{missbox}
\end{column}
\begin{column}{0.49\textwidth}
\begin{missbox}[fontupper=\footnotesize]
”Hitta alla turer och lägg dem efter varandra.” Tur 2 börjar i 3 och tur 3 i 1. Varje tur
måste skarvas in där den börjar, annars blir det ingen sammanhängande promenad.
\end{missbox}
\begin{missbox}[fontupper=\footnotesize]
”Linjär tid: $O(|V|)$.” Svaret har $|E|+1$ hörn. Linjärt betyder $O(|V|+|E|)$ — och det kräver
en pekare per hörn.
\end{missbox}
\end{column}
\end{columns}
\varstrip{Kontrast}{uppgiften}{lösningsidén, rätt och nästan rätt}{fastna, skarva, märka, räkna}
\end{frame}
